% 章节 02：状态转移模型与离散线性最优滤波递推解
\section{随机过程模型}\label{sec:model}

在处理物理现象时，仅给出经验性描述是不够的，还必须对其内在成因有所认识。若不能在某种意义上区分原因与结果——换言之，若缺乏因果性假设——便很难指望获得真正有用的结果。

% <!-- p.5 -->

人们相当普遍地接受这样一个事实：随机现象的初级宏观源是相互独立的正态随机过程。\htdfnmark{5} 一个熟知的例子是电阻器中由热骚动产生的噪声电压。在大多数情形下，\emph{被观测到}的随机现象并不能由相互独立的随机变量来描述。不同时刻观测到的随机信号之间的统计相依性（相关性），通常可解释为初级随机源与观测者之间存在一个动态系统。\emph{因此，时间的随机函数可视为由独立正态随机过程激励动态系统所产生的输出。}

正态随机信号有一条重要性质：它们通过线性系统之后仍保持为正态随机信号（\hyperref[sec:appendix]{附录}定理~\ref{thm:5} (\emph{A})）。假定初级随机源为独立正态过程，若被观测到的随机信号亦为正态过程，我们便可设观测者与初级随机源之间的动态系统是\emph{线性}的。当我们对被观测随机信号的统计性质缺乏详尽了解时，客观条件亦会迫使我们采纳这一结论：给定任何已知一阶与二阶统计矩的随机过程，总能找到一个具有完全相同一阶与二阶矩的正态随机过程（\hyperref[sec:appendix]{附录}定理~\ref{thm:5} (\emph{C})）。因此，正态分布与线性动力学是自然且相辅相成的合理假设，在统计数据匮乏时尤为如此。

如何描述一个（线性或非线性）动态系统？最根本的概念是\emph{状态}。直观而言，所谓状态，是指为了预测系统未来的行为，关于系统过去行为所必须知晓的最少定量信息（一组数值，一个函数等）。系统的动态特性随后通过\emph{状态转移}加以描述，即必须指明随着时间推移，一个状态如何转化为另一个状态。

一般而言，\emph{线性}动态系统可由如下向量微分方程与输出方程描述：
\begin{equation}\label{eq:12}
\begin{aligned}
&\frac{\mathrm{d}\bm{x}(t)}{\mathrm{d}t} = \bm{F}(t)\bm{x}(t) + \bm{D}(t)\bm{u}(t) \\
&\bm{y}(t) = \bm{M}(t)\bm{x}(t)
\end{aligned}
\end{equation}
其中 $\bm{x}(t)$ 为 $n$ 维列向量，称为系统的\emph{状态}（$\bm{x}(t)$ 的各分量 $x_i(t)$ 称为\emph{状态变量}）；$\bm{u}(t)$ 为 $m$ 维列向量（$m \leq n$），表示系统的\emph{输入}；$\bm{F}(t)$ 与 $\bm{D}(t)$ 分别为 $n \times n$ 与 $n \times m$ 矩阵。若 $\bm{F}(t)$，$\bm{D}(t)$，$\bm{M}(t)$ 的所有系数均为常数，则称动态系统 \eqref{eq:12} 是\emph{时不变}的或\emph{定常}的。最后，$\bm{y}(t)$ 为表示系统输出的 $p$ 维列向量；$\bm{M}(t)$ 为 $p \times n$ 矩阵（已校订原文维数笔误，见\hyperref[sec:k06]{校勘记 K-06}），且 $p \leq n$。

式~\eqref{eq:12} 的物理解释已在别处详加讨论 \cite{r18,r20,r23}。参看图~\ref{fig:1} 的框图会很有帮助。这不是普通的标量框图，而是一幅\emph{矩阵}框图（由表示向量信号流的粗实线标示）。图~\ref{fig:1} 中的积分器实际上代表 $n$ 个并列积分器，使得每个积分器的输出各为一个状态变量；$\bm{F}(t)$ 表示各积分器的输出如何反馈至各积分器的输入。具体而言，$f_{ij}(t)$ 即为第 $j$ 个积分器的输出反馈至第 $i$ 个积分器输入端的系数。将这一形式体系与更传统的线性系统分析方法联系起来并非难事。

\htdfig{figs/fig1.tex}{fig:1}{一般线性连续动态系统的矩阵框图}

若假定系统 \eqref{eq:12} 是定常的，且输入 $\bm{u}(t)$ 在每个采样周期内保持恒定，即满足
\begin{equation}\label{eq:13}
\bm{u}(t + \tau) = \bm{u}(t),\; 0 \leq \tau < 1,\; t = 0, 1, \dots
\end{equation}
则式~\eqref{eq:12} 可立即转化为更为便利的离散形式（对 $t = 0, 1, \dots$）：
\begin{equation*}
\bm{x}(t + 1) = \bm{\Phi}(1)\bm{x}(t) + \bm{\Delta}(1)\bm{u}(t)
\end{equation*}
其中 \cite{r18,r20}（约定 $\bm{F}^0 = \bm{I}$ 为单位矩阵）：
\begin{equation*}
\bm{\Phi}(1) = \exp(\bm{F}) = \sum_{i=0}^{\infty} \frac{\bm{F}^i}{i!}
\end{equation*}
以及
\begin{equation*}
\bm{\Delta}(1) = \left(\int_0^1 \exp(\bm{F}\tau)\,\mathrm{d}\tau\right) \bm{D}
\end{equation*}

\htdfig{figs/fig2.tex}{fig:2}{一般线性离散动态系统的矩阵框图}

参见图~\ref{fig:2}。人们亦可使用单边指数积分变换法（原称 Laplace transform methods）将 $\exp(\bm{F}\tau)$ 写成闭式形式 \cite{r18,r20,r22,r24}。若 $\bm{u}(t)$ 满足式~\eqref{eq:13} 而系统 \eqref{eq:12} 是非定常的，我们仍可类似地对 $t = 0, 1, \dots$ 写出
\begin{equation}\label{eq:14}
\begin{aligned}
&\bm{x}(t + 1) = \bm{\Phi}(t + 1; t)\bm{x}(t) + \bm{\Delta}(t)\bm{u}(t) \\
&\bm{y}(t) = \bm{M}(t)\bm{x}(t)
\end{aligned}
\end{equation}
但当然此时 $\bm{\Phi}(t + 1; t)$ 与 $\bm{\Delta}(t)$ 一般无法表为闭式形式。形如式~\eqref{eq:14} 的方程在复杂采样数据系统的研究中也经常遇到 \cite{r22}；参见图~\ref{fig:2}。

矩阵 $\bm{\Phi}(t + 1; t)$ 称为系统 \eqref{eq:12} 或 \eqref{eq:14} 的\emph{状态转移矩阵}。记号 $\bm{\Phi}(t_2; t_1)$（其中 $t_2, t_1$ 为整数）表示从时刻 $t_1$ 到时刻 $t_2$ 的状态转移。显然 $\bm{\Phi}(t; t) = \bm{I}$（单位矩阵）。若系统 \eqref{eq:12} 是定常的，则 $\bm{\Phi}(t + 1; t) = \bm{\Phi}(t + 1 - t) = \bm{\Phi}(1) = \text{常矩阵}$。同时注意乘积法则：对任意整数 $t, s, r$ 有 $\bm{\Phi}(t; s)\bm{\Phi}(s; r) = \bm{\Phi}(t; r)$；以及当转移矩阵可逆时的逆矩阵法则：$\bm{\Phi}^{-1}(t; s) = \bm{\Phi}(s; t)$。在定常系统中，$\bm{\Phi}(t; \tau) = \exp\bigl(\bm{F}(t - \tau)\bigr)$。

基于上述讨论，我们将用如下模型表示随机现象：
\begin{equation}\label{eq:15}
\bm{x}(t + 1) = \bm{\Phi}(t + 1; t)\bm{x}(t) + \bm{u}(t)
\end{equation}
其中 $\{\bm{u}(t)\}$ 为零均值向量值独立正态随机过程；鉴于\hyperref[sec:appendix]{附录}定理~\ref{thm:5} (\emph{C})，其统计特性由如下关系完全确定（已统一协方差矩阵记号为 $\bm{Q}(t)$，见\hyperref[sec:k07]{校勘记 K-07}）：
\begin{equation*}
\begin{aligned}
&\mathbb{E}[\bm{u}(t)] = \bm{0},\; \forall t \\
&\mathbb{E}[\bm{u}(t)\bm{u}^{\mathrm{T}}(s)] = \bm{0},\; t \neq s \\
&\mathbb{E}[\bm{u}(t)\bm{u}^{\mathrm{T}}(t)] = \bm{Q}(t)
\end{aligned}
\end{equation*}
% <!-- p.6 -->

当然，由\hyperref[sec:appendix]{附录}定理~\ref{thm:5} (\emph{A})，$\{\bm{x}(t)\}$ 随之也是一个零均值正态随机过程，但各时刻之间不再相互独立。事实上，若我们在稳态下考察式~\eqref{eq:15}（假定其为稳定系统），换言之，若略去初始状态 $\bm{x}(t_0)$ 的影响（令 $t_0 \to -\infty$），则有
\begin{equation*}
\bm{x}(t) = \sum_{r=-\infty}^{t-1} \bm{\Phi}(t; r + 1)\bm{u}(r)
\end{equation*}
因此当 $t \geq s$ 时，有
\begin{equation*}
\mathbb{E}[\bm{x}(t)\bm{x}^{\mathrm{T}}(s)] = \sum_{r=-\infty}^{s-1} \bm{\Phi}(t; r + 1)\bm{Q}(r)\bm{\Phi}^{\mathrm{T}}(s; r + 1)
\end{equation*}

由此可见，若我们假定一个线性动态模型并知晓正态随机激励的统计性质，便不难求出相应正态随机过程 $\{\bm{x}(t)\}$ 的统计性质。

然而在现实问题中，情况往往恰恰相反。人们给定的是协方差矩阵 $\mathbb{E}[\bm{x}(t)\bm{x}^{\mathrm{T}}(s)]$（更准确地说，是试图从有限的统计数据中估计该矩阵），而需要求解的是模型 \eqref{eq:15} 以及 $\bm{u}(t)$ 的统计性质。这是实验设计与数据归约中一个微妙且目前远未解决的难题。与有关线性最优估计问题的绝大多数工程文献一样，本文为方便起见直接从模型 \eqref{eq:15} 出发，而将如何获取模型本身视为另一个独立问题。诚然，若有可能，这两个问题\emph{理应}联合优化；但据作者所知，目前尚无任何关于这一\emph{联合}优化问题的研究。

综上所述，本文对随机过程作出如下基本假设：

\emph{物理随机现象可视为初级随机源激励动态系统所致。假定初级随机源为零均值独立正态随机过程，且动态系统为线性系统。随机过程因而由形如式~\eqref{eq:15} 的模型描述。至于表征该模型的具体数值参数如何从实验数据中获得，本文不予讨论。}

\section{线性最优估计问题的解}\label{sec:sol}

现在定义本文的核心问题。

% 问题 I：一般时变线性系统的状态估计问题定义
\begin{problemn}{I}\label{prob:I}

考虑动态模型
\begin{equation}\label{eq:16}
\bm{x}(t + 1) = \bm{\Phi}(t + 1; t)\bm{x}(t) + \bm{u}(t)
\end{equation}
\begin{equation}\label{eq:17}
\bm{y}(t) = \bm{M}(t)\bm{x}(t)
\end{equation}
其中 $\{\bm{u}(t)\}$ 是由 $n$ 维列向量构成的零均值独立正态随机过程，$\bm{x}(t)$ 为 $n$ 维状态向量，$\bm{y}(t)$ 为 $p$ 维观测向量（$p \leq n$），$\bm{\Phi}(t + 1; t)$ 与 $\bm{M}(t)$ 分别为元素是时间确定性（非随机）函数的 $n \times n$ 与 $p \times n$ 矩阵。

给定观测值 $\bm{y}(t_0), \dots, \bm{y}(t)$，试求 $\bm{x}(t_1)$ 的估计量 $\bm{x}^{\text{✫}}(t_1 \mid t)$ 使其最小化期望损失。（参看图~\ref{fig:2}，其中取 $\bm{\Delta}(t) = \bm{I}$。）

\end{problemn}

这一问题将前文提及的滤波，预测与数据平滑问题作为特例包含在内，同时也涵盖了由部分状态变量的含噪观测重构线性动态系统全部状态变量的问题（注意 $p < n$！）。

由定理~\ref{thm:2} 及评注~\ref{rem:f}（已校订引用编号，见\hyperref[sec:k08]{校勘记 K-08}）可知，问题~\ref{prob:I} 的解就是将 $\bm{x}(t_1)$ 向被观测随机变量张成的线性子空间 $\mathcal{Y}(t)$ 作正交投影 $\mathscr{P}_{\mathcal{Y}(t)}[\bm{x}(t_1)]$。如\hyperref[sec:intro]{引言}所述，这需要借助一个形如式~\eqref{eq:14} 的一般线性（未必定常！）动态系统来实现。本着这一目标，我们按如下步骤推导。

假定 $\bm{y}(t_0), \dots, \bm{y}(t - 1)$ 已经测得，即子空间 $\mathcal{Y}(t - 1)$ 已知。接着在时刻 $t$，随机向量 $\bm{y}(t)$ 被测量。同前，令 $\tilde{\bm{y}}(t \mid t - 1) = \mathscr{P}_{\mathcal{Y}(t - 1)}^{\perp}[\bm{y}(t)]$ 为 $\bm{y}(t)$ 正交于 $\mathcal{Y}(t - 1)$ 的分量。若 $\tilde{\bm{y}}(t \mid t - 1) \equiv \bm{0}$（意即几乎对每一个可能事件，该随机向量的全部坐标分量取值皆为零），则显然 $\mathcal{Y}(t)$ 与 $\mathcal{Y}(t - 1)$ 完全相同，从而测量 $\bm{y}(t)$ 不传递任何新增信息；这在具有实际物理意义的情形下通常不会发生。无论如何，$\tilde{\bm{y}}(t \mid t - 1)$ 的各坐标分量张成一个线性子空间（退化时可能为平凡零子空间 $\{0\}$），我们将其记作 $\mathcal{Z}(t)$。由定义知，$\mathcal{Y}(t - 1) \leqslant \mathcal{Y}(t)$，$\mathcal{Z}(t) \leqslant \mathcal{Y}(t)$，二者合并张成整个子空间 $\mathcal{Y}(t) = \mathcal{Y}(t - 1) \oplus \mathcal{Z}(t)$，且 $\mathcal{Z}(t)$ 中的每一个向量均正交于 $\mathcal{Y}(t - 1)$ 中的每一个向量。

按归纳法假定 $\bm{x}^{\text{✫}}(t_1 - 1 \mid t - 1)$ 已知，我们可写出：
\begin{equation}\label{eq:18}
\begin{aligned}
&\bm{x}^{\text{✫}}(t_1 \mid t) = \mathscr{P}_{\mathcal{Y}(t)}[\bm{x}(t_1)] = \mathscr{P}_{\mathcal{Y}(t - 1)}[\bm{x}(t_1)] + \mathscr{P}_{\mathcal{Z}(t)}[\bm{x}(t_1)] \\
= &\bm{\Phi}(t_1; t_1 - 1)\bm{x}^{\text{✫}}(t_1 - 1 \mid t - 1) + \mathscr{P}_{\mathcal{Y}(t - 1)}[\bm{u}(t_1 - 1)] + \mathscr{P}_{\mathcal{Z}(t)}[\bm{x}(t_1)]
\end{aligned}
\end{equation}
其中第二行系在时刻 $t_1 - 1$ 处应用式~\eqref{eq:16} 导出（已校订一般 $t_1$ 下的状态转移矩阵时间指标，见\hyperref[sec:k09]{校勘记 K-09}）。

令 $t_1 = t + s$，其中 $s$ 为任意整数。若 $s \geq 0$，则 $\bm{u}(t_1 - 1)$ 独立于 $\mathcal{Y}(t - 1)$。这是因为此时 $\bm{u}(t_1 - 1) = \bm{u}(t + s - 1)$ 独立于 $\bm{u}(t - 2), \bm{u}(t - 3), \dots$（及初始状态 $\bm{x}(t_0)$），从而由式 \eqref{eq:16}–\eqref{eq:17} 知其独立于 $\bm{y}(t_0), \dots, \bm{y}(t - 1)$，故独立于 $\mathcal{Y}(t - 1)$。又因按假设对一切 $t$ 有 $\mathbb{E}[\bm{u}(t)] = \bm{0}$，故当 $s \geq 0$ 时 $\bm{u}(t_1 - 1)$ 正交于 $\mathcal{Y}(t - 1)$。因此，若 $s \geq 0$，式~\eqref{eq:18} 第二行中间项恒为零；若 $s < 0$，计算该项则会引出相当复杂的耦合。下文仅讨论 $t_1 \geq t$ 的情形。不仅如此，我们只需详细考察 $t_1 = t + 1$ 的情形即可，因为其余情形均可轻易归结为这一情形。

当 $t_1 = t + 1$ 时，式~\eqref{eq:18} 最末一项必为作用在新息随机向量 $\tilde{\bm{y}}(t \mid t - 1)$ 上的线性运算：
\begin{equation}\label{eq:19}
\mathscr{P}_{\mathcal{Z}(t)}[\bm{x}(t + 1)] = \bm{\Delta}^{\text{✫}}(t)\tilde{\bm{y}}(t \mid t - 1)
\end{equation}
其中 $\bm{\Delta}^{\text{✫}}(t)$ 为 $n \times p$ 矩阵，上标 ${}^{\text{✫}}$ 表示“最优滤波”。

观测向量 $\bm{y}(t)$ 落在 $\mathcal{Y}(t - 1)$ 内的正交投影分量为 $\bar{\bm{y}}(t \mid t - 1) = \bm{M}(t)\bm{x}^{\text{✫}}(t \mid t - 1)$。因此
\begin{equation}\label{eq:20}
\tilde{\bm{y}}(t \mid t - 1) = \bm{y}(t) - \bar{\bm{y}}(t \mid t - 1) = \bm{y}(t) - \bm{M}(t)\bm{x}^{\text{✫}}(t \mid t - 1)
\end{equation}
联立式 \eqref{eq:18}–\eqref{eq:20}（参看图~\ref{fig:3}），我们得到
\begin{equation}\label{eq:21}
\bm{x}^{\text{✫}}(t + 1 \mid t) = \bm{\Phi}^{\text{✫}}(t + 1; t)\bm{x}^{\text{✫}}(t \mid t - 1) + \bm{\Delta}^{\text{✫}}(t)\bm{y}(t)
\end{equation}
其中
\begin{equation}\label{eq:22}
\bm{\Phi}^{\text{✫}}(t + 1; t) = \bm{\Phi}(t + 1; t) - \bm{\Delta}^{\text{✫}}(t)\bm{M}(t)
\end{equation}

由此可见，最优估计由一个与式~\eqref{eq:14} 具有相同形式的线性动态系统完成。该估计器的状态为上一步预测估计，输入为可观测随机向量最新测得的数值 $\bm{y}(t)$，状态转移矩阵由式~\eqref{eq:22} 给出。注意，线性最优滤波器的物理实现仅需：(i) 随机过程的模型；(ii) 最优新息加权矩阵算子 $\bm{\Delta}^{\text{✫}}(t)$。

估计误差同样受一个线性动态系统支配。事实上，
\begin{equation}\label{eq:23}
\begin{aligned}
&\tilde{\bm{x}}(t + 1 \mid t) = \bm{x}(t + 1) - \bm{x}^{\text{✫}}(t + 1 \mid t) \\
= &\bm{\Phi}(t + 1; t)\bm{x}(t) + \bm{u}(t) - \bm{\Phi}^{\text{✫}}(t + 1; t)\bm{x}^{\text{✫}}(t \mid t - 1) - \bm{\Delta}^{\text{✫}}(t)\bm{M}(t)\bm{x}(t) \\
= &\bm{\Phi}^{\text{✫}}(t + 1; t)\tilde{\bm{x}}(t \mid t - 1) + \bm{u}(t)
\end{aligned}
\end{equation}

\htdfig{figs/fig3.tex}{fig:3}{线性最优滤波器的矩阵框图（已校订延迟器前后信号与多步预测支路接点，见\hyperref[sec:k10]{校勘记 K-10}）}

% <!-- p.7 -->

因此，$\bm{\Phi}^{\text{✫}}(t + 1; t)$ 同时也是支配估计误差演化的线性动态系统的状态转移矩阵。

由式~\eqref{eq:23} 可立即得到最优预测误差 $\tilde{\bm{x}}(t \mid t - 1)$ 的协方差矩阵 $\bm{P}^{\text{✫}}(t)$ 所满足的递推关系。注意到 $\bm{u}(t)$ 独立于 $\bm{x}(t)$ 及历史观测从而独立于 $\tilde{\bm{x}}(t \mid t - 1)$，并利用最优残差与新息的正交性，可得
\begin{equation}\label{eq:24}
\begin{aligned}
&\bm{P}^{\text{✫}}(t + 1) = \mathbb{E}\bigl[\tilde{\bm{x}}(t + 1 \mid t)\tilde{\bm{x}}^{\mathrm{T}}(t + 1 \mid t)\bigr] \\
= &\bm{\Phi}^{\text{✫}}(t + 1; t)\mathbb{E}\bigl[\tilde{\bm{x}}(t \mid t - 1)\tilde{\bm{x}}^{\mathrm{T}}(t \mid t - 1)\bigr]\bm{\Phi}^{\text{✫}}(t + 1; t)^{\mathrm{T}} + \bm{Q}(t) \\
= &\bm{\Phi}^{\text{✫}}(t + 1; t)\mathbb{E}\bigl[\tilde{\bm{x}}(t \mid t - 1)\tilde{\bm{x}}^{\mathrm{T}}(t \mid t - 1)\bigr]\bm{\Phi}^{\mathrm{T}}(t + 1; t) + \bm{Q}(t) \\
= &\bm{\Phi}^{\text{✫}}(t + 1; t)\bm{P}^{\text{✫}}(t)\bm{\Phi}^{\mathrm{T}}(t + 1; t) + \bm{Q}(t)
\end{aligned}
\end{equation}
其中 $\bm{Q}(t) = \mathbb{E}[\bm{u}(t)\bm{u}^{\mathrm{T}}(t)]$。

接下来还需导出 $\bm{\Delta}^{\text{✫}}(t)$（进而导出 $\bm{\Phi}^{\text{✫}}(t + 1; t)$）的显式公式。由于单步新息余投影分量（已规范化单步余投影记号，见\hyperref[sec:k11]{校勘记 K-11}）
\begin{equation*}
\mathscr{P}_{\mathcal{Z}(t)}^{\perp}[\bm{x}(t + 1)] = \bm{x}(t + 1) - \mathscr{P}_{\mathcal{Z}(t)}[\bm{x}(t + 1)]
\end{equation*}
正交于 $\tilde{\bm{y}}(t \mid t - 1)$，由式~\eqref{eq:19} 立得
\begin{equation*}
\begin{aligned}
&\bm{0} = \mathbb{E}\Bigl[\bigl(\bm{x}(t + 1) - \bm{\Delta}^{\text{✫}}(t)\tilde{\bm{y}}(t \mid t - 1)\bigr)\tilde{\bm{y}}^{\mathrm{T}}(t \mid t - 1)\Bigr] \\
= &\mathbb{E}\bigl[\bm{x}(t + 1)\tilde{\bm{y}}^{\mathrm{T}}(t \mid t - 1)\bigr] - \bm{\Delta}^{\text{✫}}(t)\mathbb{E}\bigl[\tilde{\bm{y}}(t \mid t - 1)\tilde{\bm{y}}^{\mathrm{T}}(t \mid t - 1)\bigr]
\end{aligned}
\end{equation*}

注意到投影分量 $\bar{\bm{x}}(t + 1 \mid t - 1) = \mathscr{P}_{\mathcal{Y}(t - 1)}[\bm{x}(t + 1)]$ 正交于 $\mathcal{Z}(t)$，结合前文关于 $\bm{P}^{\text{✫}}(t)$ 的定义以及观测方程~\eqref{eq:17}，进一步推得（已校订过程激励误带的条件自变量，见\hyperref[sec:k11]{校勘记 K-11}）：
\begin{equation*}
\begin{aligned}
&\bm{0} = \mathbb{E}\bigl[\tilde{\bm{x}}(t + 1 \mid t - 1)\tilde{\bm{y}}^{\mathrm{T}}(t \mid t - 1)\bigr] - \bm{\Delta}^{\text{✫}}(t)\bm{M}(t)\bm{P}^{\text{✫}}(t)\bm{M}^{\mathrm{T}}(t) \\
= &\mathbb{E}\Bigl[\bigl(\bm{\Phi}(t + 1; t)\tilde{\bm{x}}(t \mid t - 1) + \bm{u}(t)\bigr)\tilde{\bm{x}}^{\mathrm{T}}(t \mid t - 1)\Bigr]\bm{M}^{\mathrm{T}}(t) - \bm{\Delta}^{\text{✫}}(t)\bm{M}(t)\bm{P}^{\text{✫}}(t)\bm{M}^{\mathrm{T}}(t)
\end{aligned}
\end{equation*}

最后，由于 $\bm{u}(t)$ 独立于 $\bm{x}(t)$ 及历史观测，有
\begin{equation*}
\bm{0} = \bm{\Phi}(t + 1; t)\bm{P}^{\text{✫}}(t)\bm{M}^{\mathrm{T}}(t) - \bm{\Delta}^{\text{✫}}(t)\bm{M}(t)\bm{P}^{\text{✫}}(t)\bm{M}^{\mathrm{T}}(t)
\end{equation*}
只要在任意时刻矩阵 $\bm{M}(t)$ 的各行均线性无关——换言之，被观测的 $p$ 个标量随机变量 $y_1(t), \dots, y_p(t)$（已校订观测分量个数下标，见\hyperref[sec:k12]{校勘记 K-12}）中没有任何一个是其余分量的线性组合——那么每当 $\bm{P}^{\text{✫}}(t)$ 正定时，矩阵 $\bm{M}(t)\bm{P}^{\text{✫}}(t)\bm{M}^{\mathrm{T}}(t)$ 便必为正定矩阵从而可逆。在此条件下，我们最终求得：
\begin{equation}\label{eq:25}
\bm{\Delta}^{\text{✫}}(t) = \bm{\Phi}(t + 1; t)\bm{P}^{\text{✫}}(t)\bm{M}^{\mathrm{T}}(t)\bigl[\bm{M}(t)\bm{P}^{\text{✫}}(t)\bm{M}^{\mathrm{T}}(t)\bigr]^{-1}
\end{equation}

由于观测始于时刻 $t_0$，故 $\tilde{\bm{x}}(t_0 \mid t_0 - 1) = \bm{x}(t_0)$；为启动借助方程~\eqref{eq:24} 对 $\bm{P}^{\text{✫}}(t)$ 的递推求值，显然必须指定初始协方差矩阵 $\bm{P}^{\text{✫}}(t_0) = \mathbb{E}[\bm{x}(t_0)\bm{x}^{\mathrm{T}}(t_0)]$。假定该矩阵正定，则方程~\eqref{eq:25} 给出 $\bm{\Delta}^{\text{✫}}(t_0)$，方程~\eqref{eq:22} 给出 $\bm{\Phi}^{\text{✫}}(t_0 + 1; t_0)$，方程~\eqref{eq:24} 给出 $\bm{P}^{\text{✫}}(t_0 + 1)$，从而完成一个递推周期。若进一步有 $\bm{Q}(t)$ 正定，则所有的 $\bm{P}^{\text{✫}}(t)$ 均保持正定，推导式~\eqref{eq:25} 所需的条件在每一步均得到满足。

现在解除 $t_1 = t + 1$ 的限制。由于 $\bm{u}(t)$ 正交于 $\mathcal{Y}(t)$，我们有
\begin{equation*}
\bm{x}^{\text{✫}}(t + 1 \mid t) = \mathscr{P}_{\mathcal{Y}(t)}\bigl[\bm{\Phi}(t + 1; t)\bm{x}(t) + \bm{u}(t)\bigr] = \bm{\Phi}(t + 1; t)\bm{x}^{\text{✫}}(t \mid t)
\end{equation*}
因此，若 $\bm{\Phi}(t + 1; t)$ 存在逆矩阵 $\bm{\Phi}(t; t + 1)$（当 $\bm{\Phi}$ 是可由微分方程描述的动态系统的状态转移矩阵时，这一条件恒成立），则有
\begin{equation*}
\bm{x}^{\text{✫}}(t \mid t) = \bm{\Phi}(t; t + 1)\bm{x}^{\text{✫}}(t + 1 \mid t)
\end{equation*}

若 $t_1 = t + s \geq t + 1$（即整数 $s \geq 1$），首先反复应用式~\eqref{eq:16} 可观察到（约定 $s = 1$ 时空和为零向量；已校订求和通项状态转移矩阵右时间指标，见\hyperref[sec:k13]{校勘记 K-13}）：
\begin{equation*}
\bm{x}(t + s) = \bm{\Phi}(t + s; t + 1)\bm{x}(t + 1) + \sum_{r=1}^{s-1} \bm{\Phi}(t + s; t + r + 1)\bm{u}(t + r)
\end{equation*}
由于 $\bm{u}(t + s - 1), \dots, \bm{u}(t + 1)$ 全部正交于 $\mathcal{Y}(t)$，对任意整数 $s \geq 1$ 有
\begin{equation*}
\begin{aligned}
&\bm{x}^{\text{✫}}(t + s \mid t) = \mathscr{P}_{\mathcal{Y}(t)}[\bm{x}(t + s)] \\
= &\mathscr{P}_{\mathcal{Y}(t)}\bigl[\bm{\Phi}(t + s; t + 1)\bm{x}(t + 1)\bigr] \\
= &\bm{\Phi}(t + s; t + 1)\bm{x}^{\text{✫}}(t + 1 \mid t)
\end{aligned}
\end{equation*}

若 $s < 0$，结果亦相类似，但递推生成固定滞后平滑估计 $\bm{x}^{\text{✫}}(t + s \mid t)$ 的动态系统将需要 $(1 - s)(n - p)$ 个状态变量（已校订平滑时标符号与动态系统维数表述，见\hyperref[sec:k14]{校勘记 K-14}）。

本节结果可归纳如下：

% 定理 3：离散卡尔曼滤波与多步预测主定理
\begin{theorem}[线性最优估计问题的解]\label{thm:3}

考察问题~\ref{prob:I}。给定观测值 $\bm{y}(t_0), \dots, \bm{y}(t)$ 时，状态 $\bm{x}(t + 1)$ 的最优一步预测估计 $\bm{x}^{\text{✫}}(t + 1 \mid t)$ 由如下线性动态系统生成：
\begin{equation*}
\bm{x}^{\text{✫}}(t + 1 \mid t) = \bm{\Phi}^{\text{✫}}(t + 1; t)\bm{x}^{\text{✫}}(t \mid t - 1) + \bm{\Delta}^{\text{✫}}(t)\bm{y}(t) \tag{21}
\end{equation*}
估计误差由如下线性动态系统给出：
\begin{equation*}
\tilde{\bm{x}}(t + 1 \mid t) = \bm{\Phi}^{\text{✫}}(t + 1; t)\tilde{\bm{x}}(t \mid t - 1) + \bm{u}(t) \tag{23}
\end{equation*}
一步预测估计误差的协方差矩阵为
\begin{equation}\label{eq:26}
\operatorname{cov}\bigl(\tilde{\bm{x}}(t \mid t - 1)\bigr) = \mathbb{E}\bigl[\tilde{\bm{x}}(t \mid t - 1)\tilde{\bm{x}}^{\mathrm{T}}(t \mid t - 1)\bigr] = \bm{P}^{\text{✫}}(t)
\end{equation}
期望二次损失为
\begin{equation}\label{eq:27}
\sum_{i=1}^{n} \mathbb{E}\bigl[\tilde{x}_i^2(t \mid t - 1)\bigr] = \operatorname{tr}\bigl(\bm{P}^{\text{✫}}(t)\bigr)
\end{equation}
当 $t \geq t_0$ 时，矩阵 $\bm{\Delta}^{\text{✫}}(t)$，$\bm{\Phi}^{\text{✫}}(t + 1; t)$ 与 $\bm{P}^{\text{✫}}(t)$ 由如下递推关系生成：
\begin{equation}\label{eq:28}
\bm{\Delta}^{\text{✫}}(t) = \bm{\Phi}(t + 1; t)\bm{P}^{\text{✫}}(t)\bm{M}^{\mathrm{T}}(t)\bigl[\bm{M}(t)\bm{P}^{\text{✫}}(t)\bm{M}^{\mathrm{T}}(t)\bigr]^{-1}
\end{equation}
\begin{equation}\label{eq:29}
\bm{\Phi}^{\text{✫}}(t + 1; t) = \bm{\Phi}(t + 1; t) - \bm{\Delta}^{\text{✫}}(t)\bm{M}(t)
\end{equation}
\begin{equation}\label{eq:30}
\bm{P}^{\text{✫}}(t + 1) = \bm{\Phi}^{\text{✫}}(t + 1; t)\bm{P}^{\text{✫}}(t)\bm{\Phi}^{\mathrm{T}}(t + 1; t) + \bm{Q}(t)
\end{equation}
为执行上述迭代计算，必须预先给定初始状态 $\bm{x}(t_0)$ 的协方差矩阵 $\bm{P}^{\text{✫}}(t_0)$ 以及随机激励 $\bm{u}(t)$ 的协方差矩阵 $\bm{Q}(t)$。最后，对任意整数 $s \geq 0$，若状态转移矩阵 $\bm{\Phi}$ 可逆，则有
\begin{equation}\label{eq:31}
\begin{aligned}
&\bm{x}^{\text{✫}}(t + s \mid t) = \bm{\Phi}(t + s; t + 1)\bm{x}^{\text{✫}}(t + 1 \mid t) \\
= &\bm{\Phi}(t + s; t + 1)\bm{\Phi}^{\text{✫}}(t + 1; t)\bm{\Phi}(t; t + s - 1)\bm{x}^{\text{✫}}(t + s - 1 \mid t - 1) \\
&\hspace{1.5em} + \bm{\Phi}(t + s; t + 1)\bm{\Delta}^{\text{✫}}(t)\bm{y}(t)
\end{aligned}
\end{equation}
因此，最优估计 $\bm{x}^{\text{✫}}(t + s \mid t)$（$s \geq 0$）同样由形如式~\eqref{eq:21} 的线性动态系统给出。

\end{theorem}

\begin{remarkn}{(h)}\label{rem:h}

当 $t \geq t_0$ 时，从式 \eqref{eq:28}–\eqref{eq:30} 中消去 $\bm{\Delta}^{\text{✫}}$ 与 $\bm{\Phi}^{\text{✫}}$，即可得到关于 $\bm{P}^{\text{✫}}(t)$ 的非线性差分方程（已校订矩阵乘积顺序，见\hyperref[sec:k15]{校勘记 K-15}）：
\begin{equation}\label{eq:32}
\begin{aligned}
&\bm{P}^{\text{✫}}(t + 1) = \bm{\Phi}(t + 1; t)\bigl\{\bm{P}^{\text{✫}}(t) - \bm{P}^{\text{✫}}(t)\bm{M}^{\mathrm{T}}(t)\bigl[\bm{M}(t)\bm{P}^{\text{✫}}(t)\bm{M}^{\mathrm{T}}(t)\bigr]^{-1}\bm{M}(t)\bm{P}^{\text{✫}}(t)\bigr\} \\
&\hspace{1.5em} \times \bm{\Phi}^{\mathrm{T}}(t + 1; t) + \bm{Q}(t)
\end{aligned}
\end{equation}
该方程仅当 $\bm{M}(t)$ 可逆时退化为线性方程，但此时随机向量 $\bm{x}(t)$ 的全部分量均可直接观测，问题归于平凡，且有 $\bm{P}^{\text{✫}}(t + 1) = \bm{Q}(t)$。注意，式~\eqref{eq:32} 在本文理论中所起的作用，恰与传统理论中的因果协方差积分方程相当。

一旦自 $t = t_0$ 起由式~\eqref{eq:32} 递推算出 $\bm{P}^{\text{✫}}(t)$，即可立即由公式 \eqref{eq:28}–\eqref{eq:29}（已校订公式引用编号，见\hyperref[sec:k16]{校勘记 K-16}）显式确定线性最优滤波器的全部系数。当然，求解方程~\eqref{eq:32}（或其对应的微分方程形式）远比求解因果协方差积分方程简便得多。

\end{remarkn}

\begin{remarkn}{(i)}\label{rem:i}

定理~\ref{thm:3} 所表述的结果尚未完全解决问题~\ref{prob:I}。例如，对于导出式~\eqref{eq:25} 所需各项假设的物理意义，非线性差分方程~\eqref{eq:32} 的收敛性与稳定性，以及最优滤波器 \eqref{eq:21} 的稳定性等问题，上文均着墨不多。这些问题实际上均能得到圆满解决，但须留待后续论文详述。在这方面，\hyperref[sec:dual]{下一节}简要提及的\hyperref[thm:4]{对偶定理}将成为核心的指导原则与分析工具，参见 \cite{r29}。

\end{remarkn}

% <!-- p.8 -->

\begin{remarkn}{(j)}\label{rem:j}

令采样周期（前文已归一化为单位 1）趋于零，即可利用本方法导出最优滤波器的微分方程表征。欲实施这一极限过渡，即从离散方程~\eqref{eq:14} 过渡到连续方程~\eqref{eq:12}，需要求取矩阵 $\bm{\Phi}^{\text{✫}}$ 的矩阵对数 $\bm{F}^{\text{✫}}$。然而，这一运算仅在 $\bm{\Phi}^{\text{✫}}$ 非奇异时方能直接施行，而不难看出事实\emph{并非}如此。其原因在于，最优滤波器实际上只需具备 $n - p$ 个状态变量，而非方程~\eqref{eq:22} 形式上所暗示的 $n$ 个状态变量。因此，经适当变换后，可将含式~\eqref{eq:22} 的滤波器方程组约化为一组等价的仅含 $n - p$ 个状态变量的方程组，且其状态转移矩阵是非奇异的。此类技术细节将在后续论著中专门讨论。

\end{remarkn}

\begin{remarkn}{(k)}\label{rem:k}

动态系统 \eqref{eq:21} 一般是非定常（时变）的，这归因于以下两点：(1) $\bm{\Phi}(t + 1; t)$ 与 $\bm{M}(t)$ 对时间 $t$ 的依赖性；(2) 估计始于有限时刻 $t = t_0$，并随着观测数据的不断积累而逐步改善。若 $\bm{\Phi}$ 与 $\bm{M}$ 均为常值矩阵，则可证明当 $t \to \infty$ 时，系统 \eqref{eq:21} 在极限意义下趋于定常动态系统；这正是经典线性估计理论所处理的情形。

\end{remarkn}

\begin{remarkn}{(l)}\label{rem:l}

值得关注的是，本文给出的全部推导既不受 $\bm{x}(t)$ 模型非平稳性的影响，也不受可用观测数据有限长度的限制。事实上，据作者所知，此前针对增长记忆滤波器给出的唯一显式递推关系出自 Blum \cite{r12}，但其所得形式远比本文结果繁复。

\end{remarkn}

\begin{remarkn}{(m)}\label{rem:m}

审视图~\ref{fig:3} 可见，最优滤波器是一个反馈系统；由于新息分量 $\tilde{\bm{y}}(t \mid t - 1)$ 显然是正交随机过程，第一求和器之后的信号实质上是白噪声。这与线性最小均方滤波中的若干熟知结论相对应，例如参见 Smith \cite{r28} 第 6 章图 6-4。然而，本文似首次从数学上给出了\emph{严格}证明：每一个线性最小均方滤波器均可通过反馈系统予以实现。不仅如此，在另一篇论文中还将证明，只要模型 \eqref{eq:16}–\eqref{eq:17} 满足相当温和的假设条件，此类反馈滤波器就必定是\emph{稳定的}，参见 \cite{r29}。

\end{remarkn}
